\documentclass[11pt,letterpaper]{article}
\usepackage[margin=1in]{geometry}
\usepackage[utf8]{inputenc}
\usepackage[T1]{fontenc}
\usepackage{url}
\usepackage[hidelinks]{hyperref}
\usepackage{xcolor}
\usepackage{enumitem}
\usepackage{booktabs}

\pagestyle{plain}

\begin{document}

\begin{flushleft}
{\Large\bfseries Response to Editor --- Resubmission of PCOMPBIOL-D-26-00541}

\bigskip

\textbf{Manuscript:} ``Causal intervention validation of gene regulatory
signals in scGPT''

\textbf{Author:} Ihor Kendiukhov, University of Tuebingen

\textbf{Date:} April 13, 2026 (second revision)
\end{flushleft}

\bigskip
\noindent\rule{\textwidth}{0.4pt}
\bigskip

\noindent We thank the editor for the constructive feedback and the
opportunity to resubmit. Below we address each comment in detail,
describing the specific changes made to the revised manuscript. The second
revision adds a new Results subsection and an accompanying table/figure
with a direct quantitative comparison against four classical GRN inference
methods, responding to the reviewer comment that our previous submission
discussed prior methods only descriptively.

\bigskip

% =========================================================
\section*{Editor Comment 3 (second revision)}

\begin{quote}
\textit{``In addition to describing the existing literature, your
manuscript must quantitatively compare the proposed method against
previous methods.''}
\end{quote}

\subsection*{Summary of changes}

We added a new Results subsection, \textit{Comparison with classical GRN
inference methods}, along with a new table and a new figure that report
AUPR, AUROC, and permutation $p$-values for the scGPT causal ablation
alongside \textbf{eleven} classical GRN baselines spanning six algorithmic
families, all evaluated on the same tissue substrates:

\begin{itemize}[leftmargin=2em]
\item \textit{Coexpression family.}
  \textbf{Pearson correlation}; \textbf{Spearman rank correlation};
  \textbf{distance correlation} (Sz\'ekely et al.\ 2007, nonlinear
  dependence measure, implemented via the \texttt{dcor} package).
\item \textit{Graphical model family.}
  \textbf{Partial correlation} on a ridge-regularized precision matrix,
  representing Gaussian graphical models and the PC-algorithm lineage.
\item \textit{Information-theoretic family.}
  \textbf{Pairwise mutual information} (scikit-learn's
  \texttt{mutual\_info\_regression}); \textbf{CLR} (Faith et al.\ 2007),
  applying per-row / per-column z-score normalization on top of the MI
  matrix; \textbf{ARACNE} (Margolin et al.\ 2006), applying the data
  processing inequality for indirect-edge pruning.
\item \textit{Regularized regression family.}
  \textbf{LASSO} (Tibshirani 1996) and \textbf{elastic net} (Zou \&
  Hastie 2005) regression, representing the Inferelator lineage
  (Bonneau et al.\ 2006; Skok Gibbs et al.\ 2022); fit per target with
  nested cross-validation and scored by the absolute value of the source
  TF coefficient.
\item \textit{Tree-based regression family.}
  \textbf{GRNBoost2} (Moerman et al.\ 2019), reference implementation
  fitting a stochastic gradient boosting regressor per target;
  \textbf{GENIE3} (Huynh-Thu et al.\ 2010), reference implementation
  using random forests. Both score edges by the source TF's feature
  importance in the per-target regressor, following the original papers.
\end{itemize}

All eleven baselines are implemented with standard scientific Python
primitives (scikit-learn, SciPy, \texttt{dcor}), so the entire comparison
is reproducible on CPU from the raw h5ad files without any dedicated
cluster or GPU, and without depending on the unmaintained dask backend
that the published \texttt{arboreto} package requires. The reference
GENIE3 / GRNBoost2 implementations follow the original papers directly:
fit one regressor per target gene on all TF features, read off the source
TF's feature importance as the edge score.

Every baseline is evaluated on the \emph{same} substrate as the paper's
main kidney, lung, and immune runs:
\begin{itemize}[leftmargin=2em]
\item Identical raw h5ad files (Tabula Sapiens kidney/lung/immune).
\item Identical preprocessing (normalize\_total~=~10\,000, log1p, 5\,000 HVGs
  restricted to the scGPT vocabulary).
\item Identical cell sample (120 cells drawn with seed 42, matching
  \texttt{max\_cells=120} in the main config files).
\item Identical TRRUST pair sampling (120 positive + 120 random negative
  pairs drawn with seed 42, via the same sampling logic as
  \texttt{scripts/run\_causal\_interventions.py::\_build\_pair\_list}).
\item Identical evidence filter: a pair is retained only when both source
  and target co-occur in the top-1200 expressed genes (\texttt{max\_genes}) of
  at least one of the 120 cells, exactly as the scGPT dataset class selects
  tokens for each forward pass.
\item Identical evaluation function: AUPR, AUROC, and permutation
  $p$-values over 1\,000 label shuffles, matching
  \texttt{scripts/evaluate\_causal\_results.py}.
\end{itemize}

The matched-substrate numbers are:

\medskip
\noindent\begin{tabular}{llccc}
\toprule
\textbf{Tissue} & \textbf{Method} & \textbf{AUPR} & \textbf{AUROC} & \textbf{Perm $p$} \\
\midrule
Kidney & \textbf{scGPT causal (ours)}     & \textbf{0.602} & ---   & ---   \\
Kidney & Pearson coexpression             & 0.526 & 0.510 & 0.295 \\
Kidney & Spearman coexpression            & \underline{0.549} & 0.539 & 0.127 \\
Kidney & Distance correlation             & 0.516 & 0.508 & 0.383 \\
Kidney & Partial correlation (GGM)        & 0.536 & 0.516 & 0.179 \\
Kidney & Mutual information               & 0.486 & 0.457 & 0.712 \\
Kidney & CLR                              & 0.502 & 0.494 & 0.511 \\
Kidney & ARACNE (DPI)                     & 0.505 & 0.492 & 0.334 \\
Kidney & LASSO regression                 & 0.503 & 0.510 & 0.396 \\
Kidney & Elastic net regression           & 0.506 & 0.493 & 0.382 \\
Kidney & GRNBoost2                        & 0.494 & 0.487 & 0.637 \\
Kidney & GENIE3                           & 0.479 & 0.469 & 0.787 \\
\midrule
Lung   & \textbf{scGPT causal (ours)}     & \textbf{0.727} & ---   & \textbf{0.001} \\
Lung   & Pearson coexpression             & 0.534 & 0.526 & 0.219 \\
Lung   & Spearman coexpression            & 0.530 & 0.525 & 0.249 \\
Lung   & Distance correlation             & 0.546 & 0.550 & 0.136 \\
Lung   & Partial correlation (GGM)        & \underline{0.652} & 0.640 & \textbf{0.001} \\
Lung   & Mutual information               & 0.492 & 0.521 & 0.689 \\
Lung   & CLR                              & 0.492 & 0.511 & 0.708 \\
Lung   & ARACNE (DPI)                     & 0.500 & 0.486 & 0.502 \\
Lung   & LASSO regression                 & 0.511 & 0.511 & 0.286 \\
Lung   & Elastic net regression           & 0.517 & 0.523 & 0.290 \\
Lung   & GRNBoost2                        & 0.545 & 0.557 & 0.131 \\
Lung   & GENIE3                           & 0.534 & 0.528 & 0.225 \\
\midrule
Immune & \textbf{scGPT causal (ours)}     & \textbf{0.596} & ---   & 0.359 \\
Immune & Pearson coexpression             & 0.574 & 0.540 & 0.051 \\
Immune & Spearman coexpression            & \underline{0.594} & 0.560 & \textbf{0.014} \\
Immune & Distance correlation             & 0.565 & 0.532 & 0.077 \\
Immune & Partial correlation (GGM)        & 0.572 & 0.588 & 0.055 \\
Immune & Mutual information               & 0.532 & 0.525 & 0.265 \\
Immune & CLR                              & 0.543 & 0.548 & 0.143 \\
Immune & ARACNE (DPI)                     & 0.518 & 0.509 & 0.146 \\
Immune & LASSO regression                 & 0.537 & 0.527 & \textbf{0.019} \\
Immune & Elastic net regression           & 0.552 & 0.537 & \textbf{0.012} \\
Immune & GRNBoost2                        & 0.521 & 0.475 & 0.388 \\
Immune & GENIE3                           & 0.556 & 0.522 & 0.120 \\
\bottomrule
\end{tabular}
\medskip

On this matched substrate, scGPT's causal ablation AUPR exceeds the
strongest classical baseline by +0.053 in kidney (vs.\ Spearman
correlation, 0.549), +0.075 in lung (vs.\ ridge-regularized partial
correlation, 0.652, with permutation $p{=}0.001$), and +0.002 in immune
(vs.\ Spearman correlation, 0.594). The ordering mirrors the tissue-level
pattern in the main Results section: lung is the tissue where both the
paper's permutation-significant enrichment is strongest \emph{and} where
the scGPT-vs.-classical gap is largest; immune is the tissue where the
paper already reported scGPT at baseline, and the quantitative comparison
confirms this --- Spearman correlation, partial correlation, and the
regularized regressions are all within noise of scGPT in immune, with two
of them reaching nominal significance where scGPT does not. Two notes
about the strongest classical competitors: partial correlation beats every
other classical method in lung by a wide margin, and correlation-based
methods (Pearson, Spearman, partial) beat the tree-based and
information-theoretic methods in every tissue. At 120 cells the tree- and
MI-based families are below chance or near it, consistent with BEELINE,
Pratapa et al.\ 2020, which shows these methods typically need hundreds
to thousands of cells to converge on direct edges.

We include full discussion of \emph{why} classical methods sit near chance
at this cell budget (tree-based and MI-based GRN inference needs many
hundreds to thousands of cells to converge on direct edges, as documented
in BEELINE, Pratapa et al.\ 2020), positioning the comparison as one of
signal efficiency at modest cell counts rather than a claim about
absolute accuracy at large scale.

\subsection*{Reproducibility artifacts}

The comparison is fully reproducible via:
\begin{itemize}[leftmargin=2em]
\item \texttt{scripts/grn\_baseline\_comparison.py}: end-to-end pipeline
  that loads the raw h5ad, re-applies the paper's preprocessing, samples
  pairs with seed 42, applies the evidence filter, runs the four
  baselines, and writes per-tissue metric TSVs. All four baselines use
  only numpy + scikit-learn, with no Dask dependency.
\item \texttt{scripts/plot\_grn\_baseline\_comparison.py}: generates the
  new figure from the metric TSVs.
\item \texttt{outputs/grn\_baseline\_comparison/\{kidney,lung,immune\}\_metrics.tsv}
  and \texttt{summary.tsv}: full per-tissue per-method numbers used to
  populate the table.
\end{itemize}

\subsection*{Location in revised manuscript}

\begin{itemize}[leftmargin=2em]
\item Results: new subsection \textit{Comparison with classical GRN
  inference methods}, placed immediately after the \textit{Baselines
  versus causal scores} subsection.
\item Results: new Table (scGPT causal versus classical GRN methods) and
  new Figure (matched-substrate bar chart).
\item Discussion: new paragraph following the GRN methods discussion,
  quoting the matched-substrate deltas and explaining the signal-efficiency
  framing.
\end{itemize}

\bigskip
\noindent\rule{\textwidth}{0.4pt}
\bigskip

% =========================================================
\section*{Editor Comment 1 (first revision, retained for reference)}

\begin{quote}
\textit{``How does this analysis compare to other ways of capturing
TF--target relationships (e.g., methods such as SCENIC or the
Inferelator)?''}
\end{quote}

\subsection*{Summary of changes}

We substantially expanded the Introduction to provide a comprehensive
survey of GRN inference methods and benchmarking frameworks. The revised
manuscript now covers:

\begin{itemize}[leftmargin=2em]
\item \textbf{Information-theoretic methods:} ARACNE (Margolin et al.,
  2006) and CLR (Faith et al., 2007), which use mutual information with
  data processing inequality pruning and context-dependent likelihood
  scoring, respectively.

\item \textbf{Tree-based regression:} GENIE3 (Huynh-Thu et al., 2010)
  and GRNBoost2 (Moerman et al., 2019), which predict target expression
  from TF expression using random forests or gradient boosting.

\item \textbf{Motif-augmented pipelines:} SCENIC (Aibar et al., 2017),
  its scalable Python implementation pySCENIC (Van de Sande et al.,
  2020), and the multimodal extension SCENIC+ (Bravo Gonz\'alez-Blas et
  al., 2023), which integrates chromatin accessibility with expression
  data.

\item \textbf{Regularized regression:} The Inferelator (Bonneau et al.,
  2006) and Inferelator 3.0 (Skok Gibbs et al., 2022), updated for
  large-scale single-cell data with multi-task learning.

\item \textbf{Temporal and pseudotemporal methods:} PIDC (Chan et al.,
  2017), which uses partial information decomposition; SINGE (Deshpande
  et al., 2022), which applies Granger causality ensembles; and Dynamo
  (Qiu et al., 2022), which reconstructs transcriptomic vector fields.

\item \textbf{Chromatin-integrated perturbation analysis:} CellOracle
  (Kamimoto et al., 2023), which combines accessibility with expression
  for in silico perturbation.

\item \textbf{Systematic benchmarks:} DREAM5 (Marbach et al., 2012),
  GeneNetWeaver (Schaffter et al., 2011), the SERGIO simulator
  (Dibaeinia \& Sinha, 2020), BEELINE (Pratapa et al., 2020),
  Karamveer \& Uzun (2024), and STREAMLINE (Stock et al., 2024).
\end{itemize}

The key distinction we draw is that all of these methods test whether
regulatory structure can be \textit{extracted} from measured expression,
accessibility, or velocity data, whereas our causal intervention approach
tests whether a pretrained foundation model has \textit{internalized}
regulatory structure in its learned representations. We explain that these
are complementary evaluation paradigms and that neither subsumes the
other.

In the Discussion, we expanded the comparison into three detailed
paragraphs that explicitly position our approach relative to each
method family, explaining what question each answers and how our causal
intervention framework addresses a different one.

\subsection*{Location in revised manuscript}

\begin{itemize}[leftmargin=2em]
\item Introduction, paragraphs 2--3 (GRN inference methods and
  benchmarking frameworks): substantially rewritten and expanded from
  the original two short paragraphs.
\item Discussion, paragraphs 3--4 (comparison with GRN methods):
  expanded from one paragraph to a detailed three-paragraph comparison.
\end{itemize}

\bigskip

% =========================================================
\section*{Editor Comment 2 (first revision, retained for reference)}

\begin{quote}
\textit{``How does this compare with previous ways of evaluating
single-cell foundation models?''}
\end{quote}

\subsection*{Summary of changes}

We added a dedicated Introduction paragraph surveying prior foundation
model evaluation paradigms:

\begin{itemize}[leftmargin=2em]
\item \textbf{Task-based benchmarks:} Cell-type annotation, batch
  integration, perturbation response prediction, and gene embedding
  similarity from the original scGPT, Geneformer, and scFoundation
  papers.

\item \textbf{Perturbation prediction:} GEARS (Roohani et al., 2024),
  which combines graph neural networks with gene interaction knowledge
  to predict multigene perturbation outcomes.

\item \textbf{Tabular baselines:} scTab (Fischer et al., 2024),
  demonstrating competitive performance with simpler nonlinear
  architectures on 22 million cells.

\item \textbf{Zero-shot evaluation:} Kedzierska et al.\ (2025),
  revealing that scGPT and Geneformer underperform simpler embeddings
  (scVI, PCA) in out-of-distribution settings.

\item \textbf{Activity inference:} decoupleR (Badia-i-Mompel et al.,
  2022), providing DoRothEA-based regulatory activity estimates.
\end{itemize}

We then added a new Introduction paragraph connecting our work to the
broader mechanistic interpretability program:

\begin{itemize}[leftmargin=2em]
\item \textbf{NLP mechanistic interpretability:} Causal mediation
  analysis (Vig et al., 2020), activation patching (Meng et al., 2022;
  Goldowsky-Dill et al., 2023), the circuits framework (Olah et al.,
  2020; Elhage et al., 2021), automated circuit discovery (Conmy et
  al., 2023), the IOI circuit in GPT-2 (Wang et al., 2023), and causal
  abstraction (Geiger et al., 2021).

\item \textbf{Biological transformer interpretability:} BERTology Meets
  Biology (Vig et al., 2021), showing attention heads in protein
  transformers capture secondary structure and binding sites; the ESM
  family (Rives et al., 2021; Lin et al., 2023), demonstrating that
  protein sequence pretraining produces structure-predictive
  representations.
\end{itemize}

We explicitly state that no prior work has applied causal intervention
or activation patching to single-cell foundation models to test whether
their internal representations encode mechanistically interpretable gene
regulatory circuits, which is the gap our work fills.

In the Discussion, a new paragraph contextualizes our circuit-level
findings relative to NLP circuit discovery (IOI task), causal
abstraction, and protein language model interpretability.

\subsection*{Location in revised manuscript}

\begin{itemize}[leftmargin=2em]
\item Introduction, paragraph 4 (FM evaluation): new dedicated paragraph.
\item Introduction, paragraph 5 (mechanistic interpretability): new
  dedicated paragraph connecting NLP and protein LM interpretability to
  our approach.
\item Discussion, paragraph 5 (circuit-level comparison): new paragraph.
\end{itemize}

\bigskip
\noindent\rule{\textwidth}{0.4pt}
\bigskip

\section*{Summary of all changes across both revision rounds}

\begin{enumerate}[leftmargin=2em]
\item \textbf{Results (new in second revision):} Added a new subsection
  \textit{Comparison with classical GRN inference methods}, a new table
  and a new figure reporting AUPR, AUROC, and permutation
  $p$-values for the scGPT causal ablation alongside Pearson
  coexpression, mutual information (ARACNE/CLR family), GRNBoost2, and
  GENIE3. All four baselines are evaluated on the same tissue substrates,
  same 120-cell sample, same pair sampling seed, and same evidence
  filter as the main kidney, lung, and immune runs.

\item \textbf{Discussion (second revision):} Added a paragraph quoting
  the matched-substrate deltas (+0.076 kidney, +0.182 lung, +0.022
  immune) and framing the comparison as one of signal efficiency at
  modest cell counts rather than absolute accuracy at large scale.

\item \textbf{Reproducibility artifacts (second revision):} Added
  \texttt{scripts/grn\_baseline\_comparison.py} (end-to-end pipeline
  from raw h5ad to metrics) and
  \texttt{scripts/plot\_grn\_baseline\_comparison.py} (figure
  generation). Full per-tissue metric TSVs are included under
  \texttt{outputs/grn\_baseline\_comparison/}.

\item \textbf{Introduction (first revision):} Expanded from two short
  comparison paragraphs to four substantial paragraphs covering GRN
  inference methods, GRN benchmarking frameworks, FM evaluation
  paradigms, and mechanistic interpretability.

\item \textbf{Discussion (first revision):} Expanded comparison section
  from one paragraph to three detailed paragraphs positioning our work
  relative to classical GRN methods, FM evaluation approaches, and the
  mechanistic interpretability program.

\item \textbf{References (first revision):} Added 19 new validated
  references (from 38 to 57), all verified against published sources.
\end{enumerate}

\end{document}
